\NSPage{061}%
\begin{EESegment}{EE-TGT-P061-001}
Because of this, the same derivative estimates still hold uniformly after the cutoffs are applied. Proposition~\ref{prop:5.3} already checked the bounds on the stress-divergence coefficients. Its estimate \eqref{eq:5.25} is exactly the finite-residual hypothesis in the lemma, with \NSi{NS-F-P061-001}{\rho_N=2h(N+1)} and no extra remainder. This gives every hypothesis of Lemma~\ref{lem:5.4} on each fixed compact range of the profile variables. Applying the lemma gives \NSi{NS-F-P061-002}{u_B,p_B,\mathcal T_{\mathrm{phys}}}. Its residual \NSi{NS-F-P061-003}{\mathcal E_B=\mathcal F_{\mathrm{slow}}(u_B,p_B,\mathcal T_{\mathrm{phys}})} is bounded by every power of \NSi{NS-F-P061-004}{q}, even after taking any of the prescribed physical derivatives.
\end{EESegment}

\begin{EESegment}{EE-TGT-P061-002}
\textbf{Step 2: keep the field divergence-free, keep its support where it belongs, and keep the velocity on the mean patch.} Differentiating the cutoff produces an extra radial term, which appears explicitly in
\end{EESegment}
\NSFormula{NS-F-P061-005}{display}{5.45}
\begin{equation}
 r u^{\mathrm{cut}}_{r,n}
 =q^{2nh}\left[\chi(c_nq)V_n-\frac{2\eta}{L}(c_nq)\chi'(c_nq)F_n\right],
 \qquad
 u^{\mathrm{cut}}_{z,n}=q^{-A+2nh}\chi(c_nq)U_n.
 \tag{5.45}\label{eq:5.45}
\end{equation}
\begin{EESegment}{EE-TGT-P061-003}
The term containing \NSi{NS-F-P061-006}{\chi'(c_nq)F_n} carries the same \NSi{NS-F-P061-007}{q^{2nh}} factor, and Lemma~\ref{lem:5.2} gives it the same radial support as the other terms. Physical derivatives commute, so the divergence remains identically zero across every cutoff transition. Equation~\eqref{eq:5.18} makes both terms inside the radial brackets zero on the mean patch. The positive-order axial and azimuthal terms vanish there as well. This proves \eqref{eq:5.44}, including the terms created when the cutoff is applied to the Stokes streamfunction.
\end{EESegment}

\begin{EESegment}{EE-TGT-P061-004}
\textbf{Step 3: prove the estimates for the normalized velocity and stress.} The extra diagonal bounds in \eqref{eq:5.40} control all normalized coefficient derivatives and Stokes streamfunction derivatives through order \NSi{NS-F-P061-008}{n}. They also control the extra streamfunction derivative and, when \NSi{NS-F-P061-009}{n\ge 2}, every derivative of \NSi{NS-F-P061-010}{\mathcal T_n/\zeta}. The normalized velocity estimates are taken on the fixed enlarged rectangle from the statement. The weighted stress estimates are taken where \NSi{NS-F-P061-011}{X_a<X<X_b} and \NSi{NS-F-P061-012}{-1\le\eta\le 1}. When an estimate reaches the axis, the profiles are written in smooth Cartesian coordinates. The finite constants include the powers of \NSi{NS-F-P061-013}{n} and the fixed derivatives of the cutoff. In addition to the physical requirements in the lemma, the chosen cutoff gives the following bound.
\end{EESegment}
\NSFormula{NS-F-P061-014}{display}{5.46}
\begin{equation}
 \left|\mathscr D^I\!\left[\chi(c_nq)q^{2nh}a_n\right]\right|
 \le 2^{-n}q^{nh}
 \qquad \bigl(n\ge\max\{2,|I|\}\bigr),
 \tag{5.46}\label{eq:5.46}
\end{equation}
\begin{EESegment}{EE-TGT-P061-005}
Here \NSi{NS-F-P061-015}{a_n} is any one of \NSi{NS-F-P061-016}{E_n,U_n,\Pi_n,V_n/\sqrt{2X},F_n/\sqrt{2X}} on the enlarged rectangle. It may instead be a component of \NSi{NS-F-P061-017}{\mathcal T_n/\zeta} on the active radial interior when \NSi{NS-F-P061-018}{n\ge 2}. The same estimate holds for the expression inside the brackets in \eqref{eq:5.45} after multiplication by \NSi{NS-F-P061-019}{q^{2nh}}. At each order, the diagonal choice in the summation lemma also includes these finitely many extra seminorms. The analytic neighborhoods may depend on the order.
\end{EESegment}

\begin{EESegment}{EE-TGT-P061-006}
Fix a derivative order. Keep the finite initial block, whose positive-order terms are all \NSi{NS-F-P061-020}{O(q^{2h})}, and sum \eqref{eq:5.46} starting at an index that is at least two and larger than that derivative order. This proves \eqref{eq:5.42}. Under the common normalization \NSi{NS-F-P061-021}{q^A}, the positive-order radial correction is actually \NSi{NS-F-P061-022}{O(q^{3h})}. For the stress components with \NSi{NS-F-P061-023}{n\ge 2}, the same argument carries the factor \NSi{NS-F-P061-024}{\zeta}. The one term with \NSi{NS-F-P061-025}{n=1} is controlled by \eqref{eq:5.22}; each derivative costs only a finite inverse power of \NSi{NS-F-P061-026}{\delta}. This proves \eqref{eq:5.43}. The wave construction can now use the fixed positive covariance representation of \NSi{NS-F-P061-027}{\mathcal T_0}, while treating the higher-order stress through signed corrections.
\end{EESegment}

\begin{EESegment}{EE-TGT-P061-007}
\textbf{Step 4: recover the full expansion and the endpoint extensions.} Fix a formal order \NSi{NS-F-P061-028}{N} and a derivative order \NSi{NS-F-P061-029}{m}. Split the summed remainder at \NSi{NS-F-P061-030}{J=\max\{2N+2,m+1\}}. The finitely many terms with \NSi{NS-F-P061-031}{N<n<J} keep their original powers, all of which are at least \NSi{NS-F-P061-032}{q^{2h(N+1)}}. The tail obeys the geometric estimate
\end{EESegment}
\NSFormula{NS-F-P061-033}{display}{none}
\[
 \sum_{n\ge J}2^{-n}q^{nh}
 \le 2^{1-J}q^{Jh}
 \le 2^{1-J}q^{2h(N+1)}
 \qquad (0<q\le 1).
\]
\NSPage{062}%
\begin{EESegment}{EE-TGT-P062-001}
The coefficient constants for the finite block in the middle may depend on \NSi{NS-F-P062-001}{N,m}. For small enough \NSi{NS-F-P062-002}{q}, every cutoff in the finite initial block is identically one. This recovers every fixed order in the original formal expansion, even after physical differentiation causes its fixed loss in the exponent. The geometric bounds are used only to sum the series, so they leave the asymptotic coefficients unchanged. Step~1 already proved the residual claim by comparing the sum with a finite partial sum taken far enough out.
\end{EESegment}

\begin{EESegment}{EE-TGT-P062-002}
Finally, the cutoff sequence grows at least geometrically. On any compact set where \NSi{NS-F-P062-003}{q} has a positive lower bound, only finitely many positive-order terms remain. Fix a physical space-time derivative. Intersect the finitely many parameter neighborhoods for the positive-order coefficients that contribute to it. Their profiles and Stokes streamfunctions extend to the resulting neighborhood. The order-zero term has the smooth one-sided extension in \NSi{NS-F-P062-004}{\eta} supplied by the leading profile and the heat construction. Since \NSi{NS-F-P062-005}{L>0}, the coordinate map carries these extensions into the physical variables, together with all the specified derivatives. This proves the claimed extensions at \NSi{NS-F-P062-006}{\eta=\pm1}.
\end{EESegment}
\end{proof}

\begin{EESegment}{EE-TGT-P062-003}
\section{The auxiliary torus and the separation of oscillatory supports}\label{sec:6}

Fix the background velocity \NSi{NS-F-P062-007}{u_B}, the pressure \NSi{NS-F-P062-008}{p_B}, and the annular stress \NSi{NS-F-P062-009}{\mathcal T_{\mathrm{phys}}} from Proposition~\ref{prop:5.5}. The momentum residual is the divergence of that stress, with the sign used in \eqref{eq:5.41}, plus a flat remainder. The waves in Proposition~\ref{prop:7.5} will supply the leading stress. Their amplitudes change under the local shear of the background according to \eqref{eq:7.5}. Before those waves can be built, they must be confined to regions where the local estimates control the variation of the shear. Their supports must also be arranged so that different waves do not create extra quadratic interactions. At the same time, this localization must allow rapid changes in time without creating radial derivatives large enough to ruin the residual estimates.
\end{EESegment}

\begin{EESegment}{EE-TGT-P062-004}
Introduce the periodic variable \NSi{NS-F-P062-010}{Y\in\mathbb T^2=\mathbb R^2/\mathbb Z^2}, independently of the physical coordinates \NSi{NS-F-P062-011}{(r,\theta,z,t)}. Physical fields are then obtained by evaluating it at the fixed map depending on \NSi{NS-F-P062-012}{(r,t)} in \eqref{eq:6.3}. This map lets time vary rapidly while adding much smaller radial derivatives. The time variation drives the pulse equation in Section~\ref{sec:7}, and the bounds on the radial derivatives control the errors caused by localization. Oscillations whose supports overlap in the remaining variables are given disjoint auxiliary supports, so their cross products vanish after evaluation. Different harmonics of the same localized oscillation still interact.
\end{EESegment}

\begin{EESegment}{EE-TGT-P062-005}
First, derive the evaluated derivative operators in \eqref{eq:6.6}. Lemma~\ref{lem:6.1} separates the supports, and Lemma~\ref{lem:6.2} puts overlapping bands into one common representation. Section~\ref{sec:6.4} then defines the coefficient classes, and Proposition~\ref{prop:6.6} gives their rules for products and derivatives.
\end{EESegment}

\begin{EESegment}{EE-TGT-P062-006}
Every preliminary geometric choice is made on the one domain \NSi{NS-F-P062-013}{0<q<q_{\mathrm{big}}}. The number \NSi{NS-F-P062-014}{q_{\mathrm{big}}} may be made smaller when the fixed base and pulse coefficients are chosen, before the correction iteration starts.
\end{EESegment}

\begin{EESegment}{EE-TGT-P062-007}
\subsection{Dyadic charts and derivatives after phase evaluation}\label{sec:6.1}

Recall from \eqref{eq:4.1} that \NSi{NS-F-P062-015}{q=q(z,t)}, \NSi{NS-F-P062-016}{A=1/2+h}, and \NSi{NS-F-P062-017}{D=1/2-h}. Work on dyadic regions where \NSi{NS-F-P062-018}{q} is comparable to a fixed scale \NSi{NS-F-P062-019}{Q}. Rescale the radial, axial, and time coordinates by \NSi{NS-F-P062-020}{Q^{1/2},Q^D,Q}, using those three factors in that order. This keeps the concentrating annulus at a bounded size in every region. For a dyadic index \NSi{NS-F-P062-021}{\ell}, set
\end{EESegment}
\NSFormula{NS-F-P062-022}{display}{6.1}
\begin{equation}
 Q=2^{-\ell},\qquad \varepsilon=Q^h,\qquad S_*=\ell^2,
 \qquad R=\frac r{\sqrt Q},\qquad Z=\frac z{Q^D},\qquad
 T=\frac\tau Q,\qquad \tau=1-t.
 \tag{6.1}\label{eq:6.1}
\end{equation}
\NSPage{063}%
\begin{EESegment}{EE-TGT-P063-001}
In these charts, the radius \NSi{NS-F-P063-001}{R} uses the fixed band scale \NSi{NS-F-P063-002}{Q}. In the earlier profile calculations, the same letter meant \NSi{NS-F-P063-003}{\sqrt{2X}=r/\sqrt q}. The two versions are related by \NSi{NS-F-P063-004}{R_{\mathrm{chart}}=\sqrt{q/Q}\,R_{\mathrm{profile}}}. The quantities \NSi{NS-F-P063-005}{Q,\varepsilon,S_*} stay constant throughout one chart. To get chart representatives, multiply a physical velocity, pressure, and residual by \NSi{NS-F-P063-006}{Q^A,Q^{2A},Q^{2A+1/2}}, using those three factors in that order. An integer matrix acts on the torus without breaking periodicity. Its two distinct eigenvalues will make the auxiliary phase vary at different rates in the radial and time directions. Fix
\end{EESegment}
\NSFormula{NS-F-P063-007}{display}{6.2}
\begin{equation}
\begin{gathered}
 J_g=\begin{pmatrix}3&1\\1&5\end{pmatrix},\qquad
 \Lambda_g=4-\sqrt2,\qquad T_g=4+\sqrt2,\qquad b_g=\sqrt2-1,\\
 v_r=(1,-b_g),\qquad v_t=(b_g,1),\qquad
 \rho_g=\frac{\log\Lambda_g}{\log T_g},\\
 \kappa_s=10^{-5},\qquad d_r=2\bigl((1+h)\rho_g-h\kappa_s\bigr)>0.
\end{gathered}
\tag{6.2}\label{eq:6.2}
\end{equation}
\begin{EESegment}{EE-TGT-P063-002}
The vectors satisfy \NSi{NS-F-P063-008}{J_gv_r=\Lambda_gv_r} and \NSi{NS-F-P063-009}{J_gv_t=T_gv_t}, with \NSi{NS-F-P063-010}{1<\Lambda_g<T_g}. Define the physical evaluation map by
\end{EESegment}
\NSFormula{NS-F-P063-011}{display}{6.3}
\begin{equation}
 Y=v_rr^{d_r}+v_tt\pmod{\mathbb Z^2}.
 \tag{6.3}\label{eq:6.3}
\end{equation}
\begin{EESegment}{EE-TGT-P063-003}
Start with a field on the extended domain, where \NSi{NS-F-P063-012}{Y} is an independent variable, and then restrict it to this map. The chain rule says that its physical derivatives are taken by the operators
\end{EESegment}
\NSFormula{NS-F-P063-013}{display}{6.4}
\begin{equation}
\begin{aligned}
 N_{\mathrm{abs}}&=v_t\mathbin\cdot\partial_Y,
 &\mathcal L_{\mathrm{abs}}&=v_r\mathbin\cdot\partial_Y,\\
 \mathfrak t&=\partial_t+N_{\mathrm{abs}},
 &\mathfrak r&=\partial_r+d_rr^{d_r-1}\mathcal L_{\mathrm{abs}}.
\end{aligned}
\tag{6.4}\label{eq:6.4}
\end{equation}
\begin{EESegment}{EE-TGT-P063-004}
Here \NSi{NS-F-P063-014}{\partial_t} and \NSi{NS-F-P063-015}{\partial_r} on the extended domain keep \NSi{NS-F-P063-016}{Y} fixed. Before restriction, \NSi{NS-F-P063-017}{\mathfrak t} represents physical differentiation in time, while \NSi{NS-F-P063-018}{\mathfrak r} represents physical differentiation in radius.
\end{EESegment}

\begin{EESegment}{EE-TGT-P063-005}
Call \NSi{NS-F-P063-019}{Y} the \emph{absolute} auxiliary variable. For band \NSi{NS-F-P063-020}{\ell}, its \emph{band representative} is the image \NSi{NS-F-P063-021}{Y_i} under the following integer covering.
\end{EESegment}
\NSFormula{NS-F-P063-022}{display}{6.5}
\begin{equation}
 Y_i=J_g^iY\pmod{\mathbb Z^2},
 \qquad
 i=i(\ell)=\left\lfloor\log_{T_g}\frac{Q^{-1-h}}{S_*}\right\rfloor.
 \tag{6.5}\label{eq:6.5}
\end{equation}
\begin{EESegment}{EE-TGT-P063-006}
If \NSi{NS-F-P063-023}{f} is a function on the band torus, its pullback to the absolute torus is \NSi{NS-F-P063-024}{f(J_g^iY)}. A function descends to that band torus, meaning that it gives one well-defined function there, when every translation \NSi{NS-F-P063-025}{Y\mapsto Y+a} with \NSi{NS-F-P063-026}{J_g^ia\in\mathbb Z^2} leaves its value unchanged. These translations are called deck translations. Only large \NSi{NS-F-P063-027}{\ell}, for which \NSi{NS-F-P063-028}{i(\ell)\ge0}, are used. Write \NSi{NS-F-P063-029}{N_i=v_t\mathbin\cdot\partial_{Y_i}} and \NSi{NS-F-P063-030}{\mathcal L_i=v_r\mathbin\cdot\partial_{Y_i}}. The eigenvector identities then give \NSi{NS-F-P063-031}{N_{\mathrm{abs}}=T_g^iN_i} and \NSi{NS-F-P063-032}{\mathcal L_{\mathrm{abs}}=\Lambda_g^i\mathcal L_i} on band pullbacks. So the physical derivatives have the chart forms below.
\end{EESegment}
\NSFormula{NS-F-P063-033}{display}{6.6}
\begin{equation}
\begin{aligned}
 \mathfrak t_*&=Q^{1+h}\mathfrak t=-\varepsilon\partial_T+c_iN_i,
 &c_i&=T_g^iQ^{1+h}\asymp S_*^{-1},\\
 D_r&=\sqrt Q\,\mathfrak r=\partial_R+M_id_rR^{d_r-1}\mathcal L_i,
 &M_i&=\Lambda_g^iQ^{d_r/2}\asymp\varepsilon^{-\kappa_s}S_*^{-\rho_g},\\
 D_z&=\sqrt Q\,\partial_z=\varepsilon\partial_Z,
 &D_\theta&=R^{-1}\partial_\theta.
\end{aligned}
\tag{6.6}\label{eq:6.6}
\end{equation}
\begin{EESegment}{EE-TGT-P063-007}
The velocity and residual normalizations introduced after \eqref{eq:6.1} are used from this point on. In \eqref{eq:6.6}, the \emph{slow} derivative terms act on \NSi{NS-F-P063-034}{R,Z,T} while keeping the auxiliary variable \NSi{NS-F-P063-035}{Y_i} fixed. The \emph{fast} terms act on \NSi{NS-F-P063-036}{Y_i}; the chain rule along the physical phase map creates them. For example, \NSi{NS-F-P063-037}{-\varepsilon\partial_T} is the slow time term, while \NSi{NS-F-P063-038}{c_iN_i} is the fast time term. The velocity and residual normalizations produce the time factor \NSi{NS-F-P063-039}{Q^{2A+1/2}Q^{-A}=Q^{1+h}} and the viscous factor \NSi{NS-F-P063-040}{Q^{2A+1/2}Q^{-A}Q^{-1}=\varepsilon}.
\end{EESegment}

\NSPage{064}%
\begin{EESegment}{EE-TGT-P064-001}
The integer \NSi{NS-F-P064-001}{i(\ell)} in \eqref{eq:6.5} satisfies
\end{EESegment}
\NSFormula{NS-F-P064-002}{display}{none}
\[
 \frac1{T_g}\frac{Q^{-1-h}}{S_*}<T_g^{i(\ell)}\le\frac{Q^{-1-h}}{S_*},
 \qquad
 \Lambda_g^{i(\ell)}=\bigl(T_g^{i(\ell)}\bigr)^{\rho_g}.
\]
\begin{EESegment}{EE-TGT-P064-002}
These inequalities give both comparisons in \eqref{eq:6.6}. The chosen covering index makes \NSi{NS-F-P064-003}{c_i} comparable to \NSi{NS-F-P064-004}{S_*^{-1}}. The choice of \NSi{NS-F-P064-005}{d_r} leaves only the factor \NSi{NS-F-P064-006}{\varepsilon^{-\kappa_s}} in \NSi{NS-F-P064-007}{M_i}, apart from its power of \NSi{NS-F-P064-008}{S_*}. These are the two derivative scales used below.
\end{EESegment}

\begin{EESegment}{EE-TGT-P064-003}
The absolute coordinate derivatives \NSi{NS-F-P064-009}{\mathfrak r,\partial_z,\partial_\theta,\mathfrak t} commute. Their only variable coefficient is a function of \NSi{NS-F-P064-010}{r} multiplying one constant direction on the torus. Cylindrical frame coefficients such as \NSi{NS-F-P064-011}{1/R} still have to be kept when vector components are differentiated. Every correction that depends on the auxiliary variable is supported away from the axis, so \NSi{NS-F-P064-012}{r^{d_r}} does not need to be smooth at \NSi{NS-F-P064-013}{r=0}. Periodicity in the auxiliary variable does not create spatial periodicity after restriction to \eqref{eq:6.3}.
\end{EESegment}

\begin{EESegment}{EE-TGT-P064-004}
The construction also integrates along each of the two torus directions. On a Fourier mode with frequency \NSi{NS-F-P064-014}{n}, this means dividing by the scalar product of that direction with \NSi{NS-F-P064-015}{n}. Both directions satisfy
\end{EESegment}
\NSFormula{NS-F-P064-016}{display}{6.7}
\begin{equation}
 |v_r\mathbin\cdot n|\ge\frac c{1+|n|},
 \qquad
 |v_t\mathbin\cdot n|\ge\frac c{1+|n|}
 \qquad \bigl(n\in\mathbb Z^2\setminus\{0\}\bigr).
 \tag{6.7}\label{eq:6.7}
\end{equation}
\begin{EESegment}{EE-TGT-P064-005}
For the first direction, \NSi{NS-F-P064-017}{v_r\mathbin\cdot n=(n_1+n_2)-\sqrt2\,n_2}. Changing the sign in front of the square root of two gives its algebraic conjugate. Multiplying the two gives the nonzero integer \NSi{NS-F-P064-018}{(n_1+n_2)^2-2n_2^2}, while the conjugate has absolute value at most \NSi{NS-F-P064-019}{C|n|}. Every nonzero integer has absolute value at least one. Since the product is such an integer and the conjugate has the stated upper bound, the original factor has the lower bound in \eqref{eq:6.7}. The same argument handles \NSi{NS-F-P064-020}{v_t}, using \NSi{NS-F-P064-021}{(n_2-n_1)+\sqrt2\,n_1}. Dividing by either dot product costs at most one power of the frequency size, which is why the inverse directional operators on zero-mean torus functions lose only finitely many torus derivatives. These operators are used in Lemma~\ref{lem:8.2} and in the equations that correct the mean.
\end{EESegment}

\begin{EESegment}{EE-TGT-P064-006}
\subsection{Slow cutoffs and disjoint auxiliary rectangles}\label{sec:6.2}

Next, assign supports to the localized oscillations. The slow cutoffs in \eqref{eq:6.9} form a squared partition of unity. When their covariances are summed in \eqref{eq:7.30}, that partition recovers the target stress. Different local oscillations are put in disjoint auxiliary rectangles whenever their slow supports meet. Lemma~\ref{lem:6.1} proves this, and the separation removes the cross terms between those oscillations.
\end{EESegment}

\begin{EESegment}{EE-TGT-P064-007}
Choose a smooth squared partition \NSi{NS-F-P064-022}{\sum_\ell\chi_\ell(q)^2=1}, with \NSi{NS-F-P064-023}{\chi_\ell} supported where \NSi{NS-F-P064-024}{q/Q\in[1/2,2]}. In each band, choose a product squared partition \NSi{NS-F-P064-025}{\sum_a\chi_{\ell,a}(R,Z,T)^2=1}. Its mesh is \NSi{NS-F-P064-026}{S_*^{-3}} in every coordinate, and each support reaches at most one mesh length past its grid point on either side. Both partitions can be made by translating one smooth bump and dividing each translate by the square root of the sum of their squares. Write
\end{EESegment}
\NSFormula{NS-F-P064-027}{display}{none}
\[
 \mathcal C_\ell(r,z,t)=\left(\frac r{\sqrt Q},\frac z{Q^D},\frac{1-t}{Q}\right),
 \qquad
 B_{\ell,a}=\supp\chi_{\ell,a}.
\]
\begin{EESegment}{EE-TGT-P064-008}
Fix a large enough starting band \NSi{NS-F-P064-028}{\ell_0}, and make \NSi{NS-F-P064-029}{q_{\mathrm{big}}} smaller if needed so that \NSi{NS-F-P064-030}{q_{\mathrm{big}}\le2^{-\ell_0}}. In this chart, the closed active part of band \NSi{NS-F-P064-031}{\ell} where \NSi{NS-F-P064-032}{\tau\ge0} is
\end{EESegment}
\NSFormula{NS-F-P064-033}{display}{none}
\[
 \mathcal A_\ell=
 \left\{\mathcal C_\ell(r,z,t):
 \begin{array}{c}
 0<q<q_{\mathrm{big}},\quad \frac12\le q/Q\le2,\quad \tau\ge0,\\[2pt]
 X_a\le r^2/(2q)\le X_b
 \end{array}
 \right\}.
\]
\begin{EESegment}{EE-TGT-P064-009}
Select the indices, representatives, and labels by
\end{EESegment}
\NSFormula{NS-F-P064-034}{display}{6.8}
\begin{equation}
\begin{gathered}
 \mathcal I_\ell=\{a:B_{\ell,a}\cap\mathcal A_\ell\ne\varnothing\},
 \qquad x^0_{\ell,a}\in B_{\ell,a}\cap\mathcal A_\ell,\\
 \Gamma=\{(\ell,a,\sigma):\ell\ge\ell_0,\ a\in\mathcal I_\ell,\ \sigma\in\{+,-\}\}.
\end{gathered}
 \tag{6.8}\label{eq:6.8}
\end{equation}
\NSPage{065}%
\begin{EESegment}{EE-TGT-P065-001}
For \NSi{NS-F-P065-001}{\gamma=(\ell,a,\sigma)\in\Gamma}, define its slow cutoff and its closed support in the physical slow coordinates by
\end{EESegment}
\NSFormula{NS-F-P065-002}{display}{6.9}
\begin{equation}
\begin{gathered}
 \eta_\gamma(R,Z,T)=\chi_\ell(q)\chi_{\ell,a}(R,Z,T),\\
 K_\gamma=\supp_{(r,z,t)}(\eta_\gamma\circ\mathcal C_\ell).
\end{gathered}
 \tag{6.9}\label{eq:6.9}
\end{equation}
\begin{EESegment}{EE-TGT-P065-002}
The supports are taken in the coordinate domain, including its smooth extension near \NSi{NS-F-P065-003}{\tau=0}. The two signs copy the same cutoff. So the partition identity on the active shell is
\end{EESegment}
\NSFormula{NS-F-P065-004}{display}{none}
\[
 \eta_{(\ell,a,+)}=\eta_{(\ell,a,-)},
 \qquad
 \sum_\ell\sum_{a\in\mathcal I_\ell}(\eta_{(\ell,a,+)}\circ\mathcal C_\ell)^2=1.
\]
\begin{EESegment}{EE-TGT-P065-003}
The representatives and every other discrete choice are fixed before any derivatives are taken.
\end{EESegment}

\begin{EESegment}{EE-TGT-P065-004}
On these supports, \NSi{NS-F-P065-005}{R} stays in one fixed compact subinterval of \NSi{NS-F-P065-006}{(0,\infty)}, while \NSi{NS-F-P065-007}{Z,T} stay in fixed bounded intervals. These bounds may depend on the fixed profile scale \NSi{NS-F-P065-008}{\mathsf C}. Derivatives of these cutoffs, of any fixed order in \NSi{NS-F-P065-009}{(R,Z,T)}, are bounded by powers of \NSi{NS-F-P065-010}{S_*}. The same remains true on fixed small enlargements inside the coordinate domain, using the smooth extension of the base from Proposition~\ref{prop:5.5}.
\end{EESegment}

\begin{EESegment}{EE-TGT-P065-005}
Give each label a small rectangle in its band torus. Choose its coordinates along \NSi{NS-F-P065-011}{v_r,v_t}, so the pulse coordinate moves at unit speed under \NSi{NS-F-P065-012}{\mathfrak t_*}. A center \NSi{NS-F-P065-013}{c_\gamma\in\mathbb T^2} will be chosen for every label. Write \NSi{NS-F-P065-014}{\pi_j(Y)=J_g^jY\pmod{\mathbb Z^2}} for the covering in \eqref{eq:6.5}. Given \NSi{NS-F-P065-015}{r_0>0}, define the band rectangle, its enlargement, and their preimages on the absolute torus by
\end{EESegment}
\NSFormula{NS-F-P065-016}{display}{6.10}
\begin{equation}
\begin{aligned}
 \mathcal R_\gamma
 &=\{c_\gamma+\xi v_r+\eta v_t:|\xi|,|\eta|<r_0\}\pmod{\mathbb Z^2},\\
 \mathcal R_\gamma^+
 &=\{c_\gamma+\xi v_r+\eta v_t:|\xi|,|\eta|<2r_0\}\pmod{\mathbb Z^2},\\
 \mathcal R_\gamma^{\mathrm{abs}}
 &=\pi_{i(\ell)}^{-1}(\mathcal R_\gamma),
 &\mathcal R_\gamma^{\mathrm{abs},+}
 &=\pi_{i(\ell)}^{-1}(\mathcal R_\gamma^+).
\end{aligned}
\tag{6.10}\label{eq:6.10}
\end{equation}
\begin{EESegment}{EE-TGT-P065-006}
Take \NSi{NS-F-P065-017}{r_0} small enough that the enlarged parametrizations are injective, so they never give the same point from two different coordinate pairs. Choose a representative of \NSi{NS-F-P065-018}{c_\gamma} in \NSi{NS-F-P065-019}{\mathbb R^2}. On a local lift of \NSi{NS-F-P065-020}{\mathcal R_\gamma}, the coordinates are
\end{EESegment}
\NSFormula{NS-F-P065-021}{display}{6.11}
\begin{equation}
\begin{gathered}
 Y_i-c_\gamma-k_{\mathrm{copy}}=\xi_gv_r+\eta_gv_t,
 \qquad |\xi_g|,|\eta_g|<r_0,\\
 v=\frac{\eta_g+r_0}{c_i},
 \qquad
 L_s=\frac{2r_0}{c_i}\asymp S_*.
\end{gathered}
\tag{6.11}\label{eq:6.11}
\end{equation}
\begin{EESegment}{EE-TGT-P065-007}
Here \NSi{NS-F-P065-022}{k_{\mathrm{copy}}\in\mathbb Z^2} chooses the local lift, and the same coordinates extend to \NSi{NS-F-P065-023}{\mathcal R_\gamma^+}. The dual coordinates satisfy \NSi{NS-F-P065-024}{\mathcal L_i\eta_g=N_i\xi_g=0} and \NSi{NS-F-P065-025}{N_i\eta_g=\mathcal L_i\xi_g=1}. Since \NSi{NS-F-P065-026}{c_i} is constant within the chart, \eqref{eq:6.6} gives
\end{EESegment}
\NSFormula{NS-F-P065-027}{display}{6.12}
\begin{equation}
 D_rv=D_zv=0,
 \qquad \mathfrak t_*v=1,
 \qquad N_i\xi_g=0.
 \tag{6.12}\label{eq:6.12}
\end{equation}
\begin{EESegment}{EE-TGT-P065-008}
So \NSi{NS-F-P065-028}{v} measures normalized time along a rectangle and stays constant under the normalized spatial derivatives. The centers will now be chosen so that these local coordinates can be used without creating cross products between different local oscillations.
\end{EESegment}

\begin{lemma}\label{lem:6.1}
\begin{EESegment}{EE-TGT-P065-009}
For the labels \NSi{NS-F-P065-029}{\Gamma} and slow supports \NSi{NS-F-P065-030}{K_\gamma} in Equations~\eqref{eq:6.8} and \eqref{eq:6.9}, there are centers \NSi{NS-F-P065-031}{c_\gamma} and one radius \NSi{NS-F-P065-032}{r_0>0}, shared by every label and independent of the band, with two properties. First, the enlarged rectangles in \eqref{eq:6.10} are parametrized injectively. Second,
\end{EESegment}
\NSFormula{NS-F-P065-033}{display}{6.13}
\begin{equation}
 \gamma\ne\gamma',\qquad K_\gamma\cap K_{\gamma'}\ne\varnothing
 \qquad\Longrightarrow\qquad
 \overline{\mathcal R_\gamma^{\mathrm{abs},+}}
 \cap
 \overline{\mathcal R_{\gamma'}^{\mathrm{abs},+}}
 =\varnothing.
 \tag{6.13}\label{eq:6.13}
\end{equation}
\end{lemma}
\begin{proof}
\begin{EESegment}{EE-TGT-P065-010}
\textbf{Step 1: bound how many interactions are possible.} If two dyadic supports meet, then \NSi{NS-F-P065-034}{q/Q,q/Q'\in[1/2,2]}, so \NSi{NS-F-P065-035}{|\ell-\ell'|\le2}. Enlarge every mesh box by one fixed factor. Join two different labels when their enlarged slow boxes meet in physical coordinates and their band indices differ by at most four. Also join the two opposite signs of the same box. The resulting graph has
\end{EESegment}
\NSPage{066}%
\begin{EESegment}{EE-TGT-P066-001}
bounded degree. In fact, when \NSi{NS-F-P066-001}{|\ell-\ell'|\le4}, the change of physical scale in each coordinate is bounded, and so is the ratio \NSi{NS-F-P066-002}{\ell^2/(\ell')^2}. A box can then meet only a bounded number of enlarged boxes in any of those bands. Neither the degrees nor the number of colors needed depends on the band.
\end{EESegment}

\begin{EESegment}{EE-TGT-P066-002}
The differences between the covering indices \NSi{NS-F-P066-003}{i(\ell)} also have one fixed bound. If every band has \NSi{NS-F-P066-004}{\ell\ge\ell_0}, then \NSi{NS-F-P066-005}{|\ell-\ell'|\le4} gives
\end{EESegment}
\NSFormula{NS-F-P066-006}{display}{6.14}
\begin{equation}
 |i(\ell)-i(\ell')|
 \le 1+\frac{4(1+h)\log2+8/\ell_0}{\log T_g}
 \le\Delta_{\max},
 \tag{6.14}\label{eq:6.14}
\end{equation}
\begin{EESegment}{EE-TGT-P066-003}
after choosing an integer upper bound. To get this estimate, subtract \NSi{NS-F-P066-007}{((1+h)\ell\log2-2\log\ell)/\log T_g} at the two band indices and use the mean value theorem. Taking the floor adds at most one.
\end{EESegment}

\begin{EESegment}{EE-TGT-P066-004}
\textbf{Step 2: separate the centers under the coverings that matter.} Color the countable graph greedily, using only finitely many colors. Choose one rational center \NSi{NS-F-P066-008}{c_\nu\in\mathbb T^2} for each color \NSi{NS-F-P066-009}{\nu}, and set \NSi{NS-F-P066-010}{c_\gamma=c_\nu} whenever \NSi{NS-F-P066-011}{\gamma} has that color. Choose the color centers so that all the following ordered constraints hold.
\end{EESegment}
\NSFormula{NS-F-P066-012}{display}{6.15}
\begin{equation}
 c_\mu\ne J_g^\Delta c_\nu\pmod{\mathbb Z^2}
 \quad (0\le\Delta\le\Delta_{\max}),
 \qquad
 \text{unless }(\Delta,\nu,\mu)=(0,\nu,\nu).
 \tag{6.15}\label{eq:6.15}
\end{equation}
\begin{EESegment}{EE-TGT-P066-005}
Every excluded equality is a proper closed condition on the finite tuple of centers. When the two colors are the same and \NSi{NS-F-P066-013}{\Delta>0}, this follows because \NSi{NS-F-P066-014}{J_g^\Delta-I} is invertible. The complement of all these conditions is open and dense, and it contains a rational tuple. This leaves every one of the finitely many forbidden differences a positive distance from the integer lattice.
\end{EESegment}

\begin{EESegment}{EE-TGT-P066-006}
\textbf{Step 3: choose one size for every rectangle.} Take rectangles in the eigenvector coordinates around these centers. Make them small enough both to be injective on the band torus and to keep the separation after the fixed enlargement. If a point belonged to the lifted rectangles of two adjacent labels at levels \NSi{NS-F-P066-015}{i} and \NSi{NS-F-P066-016}{i+\Delta}, it would give \NSi{NS-F-P066-017}{c_\mu-J_g^\Delta c_\nu=J_g^\Delta e_\nu-e_\mu} modulo the lattice, with \NSi{NS-F-P066-018}{|e_\nu|+|e_\mu|\le Cr_0}. Since \NSi{NS-F-P066-019}{\Delta} has a fixed finite bound, one fixed \NSi{NS-F-P066-020}{r_0} small enough makes this contradict \eqref{eq:6.15}. Step~1 joined every pair whose slow supports overlap. This proves \eqref{eq:6.13}, including the separation needed for the supports of derivatives.
\end{EESegment}
\end{proof}

\begin{EESegment}{EE-TGT-P066-007}
For the rest of the construction, fix the centers \NSi{NS-F-P066-021}{c_\gamma} and the radius \NSi{NS-F-P066-022}{r_0} from Lemma~\ref{lem:6.1}, along with one integer \NSi{NS-F-P066-023}{\Delta_{\max}} that satisfies \eqref{eq:6.14}. In particular, take a family \NSi{NS-F-P066-024}{(F_\gamma)_{\gamma\in\Gamma}} whose supports in \NSi{NS-F-P066-025}{(r,z,t,Y)} are uniform in \NSi{NS-F-P066-026}{\theta}. Then
\end{EESegment}
\NSFormula{NS-F-P066-027}{display}{none}
\[
 \bigl[\supp F_\gamma\subset K_\gamma\times\mathcal R_\gamma^{\mathrm{abs},+}
 \text{ for all }\gamma\in\Gamma\bigr]
 \qquad\Longrightarrow\qquad
 F_\gamma F_{\gamma'}=0
 \quad(\gamma\ne\gamma').
\]
\begin{EESegment}{EE-TGT-P066-008}
The same identity holds for derivatives of smoothly extended fields, and it still holds after evaluation on \eqref{eq:6.3}.
\end{EESegment}

\begin{EESegment}{EE-TGT-P066-009}
Choose a nonzero smooth transverse cutoff \NSi{NS-F-P066-028}{\chi_g\in C_c^\infty((-r_0,r_0))}, with \NSi{NS-F-P066-029}{0\le\chi_g\le1}. Rescale one fixed function to get the time cutoff \NSi{NS-F-P066-030}{0\le\psi\le1}, with
\end{EESegment}
\NSFormula{NS-F-P066-031}{display}{6.16}
\begin{equation}
 \psi(v)=1\quad\text{if }|v-L_s/2|\le L_s/5,
 \qquad
 \supp\psi\subset\{|v-L_s/2|<L_s/3\}.
 \tag{6.16}\label{eq:6.16}
\end{equation}
\begin{EESegment}{EE-TGT-P066-010}
These cutoffs stay strictly inside the enlarged rectangles. Once the pulse equation has been solved there, multiplication by \NSi{NS-F-P066-032}{\psi} keeps its time support inside the rectangle.
\end{EESegment}

\begin{EESegment}{EE-TGT-P066-011}
\subsection{One common torus for overlapping bands}\label{sec:6.3}

The rectangle \NSi{NS-F-P066-033}{\mathcal R_\gamma} from \eqref{eq:6.10}, for a label in band \NSi{NS-F-P066-034}{\ell}, is written in the band coordinate \NSi{NS-F-P066-035}{Y_i=\pi_{i(\ell)}(Y)} from \eqref{eq:6.5}. Before fields from overlapping bands are added, put all these rectangles and fields on the same torus \NSi{NS-F-P066-036}{H} without changing their physical values. Lemma~\ref{lem:6.2} gives uniformly bounded derivative factors
\end{EESegment}
